% questions.tex —— 2012 年·安徽卷·理科 · 题目内容片段（无答案）
% 由 tools/generate_projects.py 依来源总账自 2012/2012_数学真题_安徽理_试题.md 生成，勿手改（改源 MD 后重跑）

\gkpart{一、选择题}

\begin{question}
\gkti{1}复数 $z$ 满足 $(z-\mathrm i)(2-\mathrm i)=5$，则 $z=$\mbox{（　　）}

（A）$-2-2\mathrm i$；（B）$-2+2\mathrm i$；（C）$2-2\mathrm i$；（D）$2+2\mathrm i$。

\end{question}
\begin{question}
\gkti{2}下列函数中，不满足 $f(2x)=2f(x)$ 的是\mbox{（　　）}

（A）$f(x)=|x|$；（B）$f(x)=x-|x|$；（C）$f(x)=x+1$；（D）$f(x)=-x$。

\Needspace{0.4\textheight}
\end{question}
\begin{question}
\gkti{3}如图所示，程序框图（算法流程图）的输出结果是\mbox{（　　）}
\nopagebreak[4]
\begin{center}% 2012 年安徽卷文科第 3 题：输出 y 的程序框图。
% 来源 PDF：普通高考/2012/2012安徽理.pdf，第 1 页（印面 715）左栏。
% 流程：开始 → x=1,y=1 → x≤4? 是：x=2x，y=y+1，回到判断；否：输出 y → 结束。
\begin{tikzpicture}[semithick, >=Stealth,
  block/.style={draw, minimum width=3.2cm, minimum height=0.62cm, align=center},
  dec/.style={draw, diamond, aspect=2.6, inner sep=1pt, align=center},
  term/.style={draw, rounded corners=6pt, minimum width=1.5cm, minimum height=0.52cm},
  io/.style={draw, trapezium, trapezium left angle=70, trapezium right angle=110,
             minimum width=2.2cm, minimum height=0.58cm},
  arr/.style={->}]
  \node (start) [term] at (0,6.6) {开始};
  \node (init)  [block] at (0,5.7) {$x=1,\ y=1$};
  \node (dec)   [dec]  at (0,4.35) {$x\leqslant4$?};
  \node (body)  [block] at (3.7,4.35) {$x=2x,\ y=y+1$};
  \node (out)   [io]   at (0,2.95) {输出 $y$};
  \node (end)   [term] at (0,1.95) {结束};
  \draw[arr] (start) -- (init);
  \draw[arr] (init) -- (dec);
  \draw[arr] (dec) -- node[above] {是} (body);
  \draw[arr] (body.north) -- (3.7,5.0) -- (0,5.0);
  \draw[arr] (dec) -- node[right] {否} (out);
  \draw[arr] (out) -- (end);
\end{tikzpicture}
\end{center}
\nopagebreak[4]
（A）$3$；（B）$4$；（C）$5$；（D）$8$。

\end{question}
\begin{question}
\gkti{4}公比为 $\sqrt[3]{2}$ 的等比数列 $\{a_n\}$ 的各项都是正数，且 $a_3a_{11}=16$，则 $\log_2a_{16}=$\mbox{（　　）}

（A）$4$；（B）$5$；（C）$6$；（D）$7$。

\Needspace{0.4\textheight}
\end{question}
\begin{question}
\gkti{5}甲、乙两人在一次射击比赛中各射靶 $5$ 次，两人成绩的条形统计图如图所示，则\mbox{（　　）}
\nopagebreak[4]
\begin{center}% 2012 年安徽卷理科第 5 题：甲、乙两人各射靶 5 次的成绩条形统计图。
% 来源 PDF：普通高考/2012/2012安徽理.pdf，第 1 页（印面 713）左栏。
% 甲：4,5,6,7,8 各 1 次；乙：5 环 3 次、6 环 1 次、9 环 1 次。
\begin{tikzpicture}[scale=0.62, semithick, font=\small]
  % ---------- 甲 ----------
  \begin{scope}
    \draw[->] (0,0) -- (8.4,0) node[below] {环数};
    \draw[->] (0,0) -- (0,4.2) node[left] {频数};
    \foreach \y in {1,2,3} \draw (0,\y) -- (-0.12,\y) node[left] {\y};
    \foreach \x/\s in {1/4,2/5,3/6,4/7,5/8} {
      \draw (\x-0.42,0) rectangle (\x+0.42,1);
      \draw (\x,-0.12) -- (\x,-0.12) node[below] {\s};
    }
    \node at (2.5,-1.55) {（甲）};
  \end{scope}
  % ---------- 乙 ----------
  \begin{scope}[shift={(11.2,0)}]
    \draw[->] (0,0) -- (8.4,0) node[below] {环数};
    \draw[->] (0,0) -- (0,4.2) node[left] {频数};
    \foreach \y in {1,2,3} \draw (0,\y) -- (-0.12,\y) node[left] {\y};
    \foreach \x/\s/\h in {1/4/0,2/5/3,3/6/1,4/7/0,5/8/0,6/9/1} {
      \ifnum\h>0 \draw (\x-0.42,0) rectangle (\x+0.42,\h); \fi
      \draw (\x,-0.12) -- (\x,-0.12) node[below] {\s};
    }
    \node at (3.0,-1.55) {（乙）};
  \end{scope}
\end{tikzpicture}
\end{center}
\nopagebreak[4]
（A）甲的成绩的平均数小于乙的成绩的平均数；（B）甲的成绩的中位数等于乙的成绩的中位数；（C）甲的成绩的方差小于乙的成绩的方差；（D）甲的成绩的极差小于乙的成绩的极差。

\end{question}
\begin{question}
\gkti{6}设平面 $\alpha$ 与平面 $\beta$ 相交于直线 $m$，直线 $a$ 在平面 $\alpha$ 内，直线 $b$ 在平面 $\beta$ 内，且 $b\perp m$，则“$\alpha\perp\beta$”是“$a\perp b$”的\mbox{（　　）}

（A）充分不必要条件；（B）必要不充分条件；（C）充分必要条件；（D）既不充分也不必要条件。

\end{question}
\begin{question}
\gkti{7}$\left(x^2+2\right)\left(\dfrac1{x^2}-1\right)^5$ 的展开式的常数项是\mbox{（　　）}

（A）$-3$；（B）$-2$；（C）$2$；（D）$3$。

\end{question}
\begin{question}
\gkti{8}在平面直角坐标系中，点 $O(0,0)$，$P(6,8)$，将向量 $\overrightarrow{OP}$ 绕点 $O$ 按逆时针方向旋转 $\dfrac{3\pi}4$ 后得向量 $\overrightarrow{OQ}$，则点 $Q$ 的坐标是\mbox{（　　）}

（A）$\left(-7\sqrt2,-\sqrt2\right)$；（B）$\left(-7\sqrt2,\sqrt2\right)$；（C）$\left(-4\sqrt6,-2\right)$；（D）$\left(-4\sqrt6,2\right)$。

\end{question}
\begin{question}
\gkti{9}过抛物线 $y^2=4x$ 的焦点 $F$ 的直线交该抛物线于 $A$，$B$ 两点，$O$ 为坐标原点。若 $|AF|=3$，则 $\triangle AOB$ 的面积为\mbox{（　　）}

（A）$\dfrac{\sqrt2}2$；（B）$\sqrt2$；（C）$\dfrac{3\sqrt2}2$；（D）$2\sqrt2$。

\end{question}
\begin{question}
\gkti{10}$6$ 位同学在毕业聚会活动中进行纪念品的交换，任意两位同学之间最多交换一次，进行交换的两位同学互赠一份纪念品。已知 $6$ 位同学之间共进行了 $13$ 次交换，则收到 $4$ 份纪念品的同学人数为\mbox{（　　）}

（A）$1$ 或 $3$；（B）$1$ 或 $4$；（C）$2$ 或 $3$；（D）$2$ 或 $4$。

\end{question}
\gkpart{二、填空题}

\begin{question}
\gkti{11}若 $x$，$y$ 满足约束条件 $\begin{cases}x\geqslant0,\\x+2y\geqslant3,\\2x+y\leqslant3,\end{cases}$ 则 $x-y$ 的取值范围是\gkblank{}。

\Needspace{0.4\textheight}
\end{question}
\begin{question}
\gkti{12}某几何体的三视图如图所示，该几何体的表面积是\gkblank{}。
\nopagebreak[4]
\begin{center}% 2012 年安徽卷文科第 12 题：直棱柱的三视图（单位：cm，图中只标数值）。
% 来源 PDF：普通高考/2012/2012安徽文.pdf，第 1 页（印面 715）中栏。
% 正视图 5×4，内部竖线在 x=2；侧视图 4×4；俯视图为直角梯形，上底 5、下底 2、高 4。
\begin{tikzpicture}[scale=0.70, semithick,
  dim/.style={<->, >=Stealth, thin},
  lab/.style={font=\small}]
  % ---------- 正视图 ----------
  \begin{scope}
    \draw (0,0) rectangle (5,4);
    \draw (2,0) -- (2,4);
    \draw[dim] (-0.45,0) -- (-0.45,4);
    \node[left] at (-0.45,2) {4};
    \node[below] at (2.5,-0.45) {正（主）视图};
  \end{scope}
  % ---------- 侧视图 ----------
  \begin{scope}[shift={(7.6,0)}]
    \draw (0,0) rectangle (4,4);
    \node[below] at (2,-0.45) {侧（左）视图};
  \end{scope}
  % ---------- 俯视图 ----------
  \begin{scope}[shift={(1.2,-6.6)}]
    \draw (0,0) -- (2,0) -- (5,4) -- (0,4) -- cycle;
    \draw[dim] (5.45,0) -- (5.45,4);
    \node[right] at (5.45,2) {4};
    \draw[dim] (0,4.4) -- (5,4.4);
    \node[above] at (2.5,4.4) {5};
    \draw[dim] (0,-0.4) -- (2,-0.4);
    \node[below] at (1,-0.4) {2};
    \node[below] at (2.5,-1.05) {俯视图};
  \end{scope}
\end{tikzpicture}
\end{center}
\nopagebreak[4]
\end{question}
\begin{question}
\gkti{13}在极坐标系中，圆 $\rho=4\sin\theta$ 的圆心到直线 $\theta=\dfrac\pi6$（$\rho\in\mathbf R$）的距离是\gkblank{}。

\end{question}
\begin{question}
\gkti{14}若平面向量 $\boldsymbol a$，$\boldsymbol b$ 满足 $|2\boldsymbol a-\boldsymbol b|\leqslant3$，则 $\boldsymbol a\cdot\boldsymbol b$ 的最小值是\gkblank{}。

\end{question}
\begin{question}
\gkti{15}设 $\triangle ABC$ 的内角 $A$，$B$，$C$ 所对边的长分别为 $a$，$b$，$c$，则下列命题正确的是\gkblank{}。（写出所有正确命题的编号）

① 若 $ab>c^2$，则 $C<\dfrac\pi3$；

② 若 $a+b>2c$，则 $C<\dfrac\pi3$；

③ 若 $a^3+b^3=c^3$，则 $C<\dfrac\pi2$；

④ 若 $(a+b)c<2ab$，则 $C>\dfrac\pi2$；

⑤ 若 $(a^2+b^2)c^2<2a^2b^2$，则 $C>\dfrac\pi2$。

\end{question}
\gkpart{三、解答题}

\begin{question}
\gkti{16}设函数 $f(x)=\dfrac{\sqrt2}2\cos\left(2x+\dfrac\pi4\right)+\sin^2x$。

（1） 求 $f(x)$ 的最小正周期；

（2） 设函数 $g(x)$ 对任意 $x\in\mathbf R$，有 $g\left(x+\dfrac\pi2\right)=g(x)$，且当 $x\in\left[0,\dfrac\pi2\right]$ 时，$g(x)=-\dfrac12f(x)$。求 $g(x)$ 在区间 $[-\pi,0]$ 上的解析式。

\end{question}
\begin{question}
\gkti{17}某单位招聘面试，每次从试题库中随机调用一道试题。若调用的是 A 类型试题，则使用后该试题回库，并增补一道 A 类型试题和一道 B 类型试题入库，此次调题工作结束；若调用的是 B 类型试题，则使用后该试题回库，此次调题工作结束。试题库中现共有 $n+m$ 道试题，其中有 $n$ 道 A 类型试题和 $m$ 道 B 类型试题。以 $X$ 表示两次调题工作完成后，试题库中 A 类试题的数量。

（1） 求 $X=n+2$ 的概率；

（2） 设 $m=n$，求 $X$ 的分布列和均值（数学期望）。

\Needspace{0.4\textheight}
\end{question}
\begin{question}
\gkti{18}平面图形 $ABB_1A_1C_1C$ 如图 1 所示，其中 $BB_1C_1C$ 是矩形，$BC=2$，$BB_1=4$，$AB=AC=\sqrt2$，$A_1B_1=A_1C_1=\sqrt5$。现将该平面图形分别沿 $BC$ 和 $B_1C_1$ 折叠，使 $\triangle ABC$ 与 $\triangle A_1B_1C_1$ 所在平面都与平面 $BB_1C_1C$ 垂直，再分别连接 $A_1A$，$A_1B$，$A_1C$，得到如图 2 所示的空间图形，对此空间图形解答下列问题。

（1） 证明：$AA_1\perp BC$；

（2） 求 $AA_1$ 的长；

（3） 求二面角 $A-BC-A_1$ 的余弦值。
\nopagebreak[4]
\begin{center}% 2012 年安徽卷理科第 18 题：平面展开图（图 1）与折叠后的空间图形（图 2）。
% 来源 PDF：普通高考/2012/2012安徽理.pdf，第 2 页（印面 714）第 1 栏。
% 图 1：矩形 BB1C1C（BC=2、BB1=4），上接等腰直角三角形 ABC，下接等腰三角形 A1B1C1。
% 图 2：A、A1 分别折到矩形所在平面的两侧，A、A1 到对应折线的距离为 1 与 2。
\begin{tikzpicture}[semithick, font=\small]
  % ---------- 图 1 ----------
  \begin{scope}
    \coordinate (B) at (0,4);
    \coordinate (C) at (2,4);
    \coordinate (B1) at (0,0);
    \coordinate (C1) at (2,0);
    \coordinate (A) at (1,5);
    \coordinate (A1) at (1,-2);
    \draw (B) -- (C) -- (C1) -- (B1) -- cycle;
    \draw (A) -- (B) -- (C) -- cycle;
    \draw (A1) -- (B1) -- (C1) -- cycle;
    \node[above] at (A) {$A$};
    \node[left] at (B) {$B$};
    \node[right] at (C) {$C$};
    \node[left] at (B1) {$B_1$};
    \node[right] at (C1) {$C_1$};
    \node[below] at (A1) {$A_1$};
    \node[below] at (1,-2.7) {图 1};
  \end{scope}
  % ---------- 图 2 ----------
  \begin{scope}[shift={(7.4,0)}]
    \coordinate (C) at (1.4,4.6);
    \coordinate (B) at (2.5,4.0);
    \coordinate (A) at (0.3,3.6);
    \coordinate (C1) at (1.2,0.9);
    \coordinate (B1) at (2.3,0.5);
    \coordinate (A1) at (4.4,0.7);
    \draw (A) -- (C) -- (B) -- cycle;
    \draw (C) -- (C1);
    \draw (B) -- (B1);
    \draw (B1) -- (C1);
    \draw (A) -- (A1);
    \draw (A1) -- (B);
    \draw (A1) -- (C);
    \draw (A1) -- (B1);
    \draw[dashed] (C1) -- (A1);
    \node[above] at (C) {$C$};
    \node[right] at (B) {$B$};
    \node[left] at (A) {$A$};
    \node[left] at (C1) {$C_1$};
    \node[below] at (B1) {$B_1$};
    \node[right] at (A1) {$A_1$};
    \node[below] at (2.2,-0.1) {图 2};
  \end{scope}
\end{tikzpicture}
\end{center}
\nopagebreak[4]
\end{question}
\begin{question}
\gkti{19}设函数 $f(x)=a\mathrm e^x+\dfrac1{a\mathrm e^x}+b$（$a>0$）。

（1） 求 $f(x)$ 在 $[0,+\infty)$ 内的最小值；

（2） 设曲线 $y=f(x)$ 在点 $(2,f(2))$ 处的切线方程为 $y=\dfrac32x$，求 $a$，$b$ 的值。

\Needspace{0.4\textheight}
\end{question}
\begin{question}
\gkti{20}如图，点 $F_1(-c,0)$，$F_2(c,0)$ 分别是椭圆 $C:\dfrac{x^2}{a^2}+\dfrac{y^2}{b^2}=1$（$a>b>0$）的左、右焦点，过点 $F_1$ 作 $x$ 轴的垂线交椭圆 $C$ 的上半部分于点 $P$，过点 $F_2$ 作直线 $PF_2$ 的垂线交直线 $x=\dfrac{a^2}c$ 于点 $Q$。

（1） 如果点 $Q$ 的坐标是 $(4,4)$，求此时椭圆 $C$ 的方程；

（2） 证明：直线 $PQ$ 与椭圆 $C$ 只有一个交点。
\nopagebreak[4]
\begin{center}% 2012 年安徽卷理科第 20 题：椭圆焦点 F1、F2，P 在 x=-c 与椭圆的上半部分交点，Q 为垂线与 x=a^2/c 的交点。
% 来源 PDF：普通高考/2012/2012安徽理.pdf，第 2 页（印面 714）第 2 栏。
% 按 (1) 的结果 a=2、c=1、b=√3 作图：P(-1,1.5)，Q(4,4)，F2(1,0)。
\begin{tikzpicture}[scale=0.95, semithick, >=Stealth, font=\small]
  \draw (0,0) ellipse (2 and 1.732);
  \draw[->] (-3.2,0) -- (4.9,0) node[right] {$x$};
  \draw[->] (0,-2.1) -- (0,4.6) node[above] {$y$};
  \coordinate (O) at (0,0);
  \coordinate (F1) at (-1,0);
  \coordinate (F2) at (1,0);
  \coordinate (P) at (-1,1.5);
  \coordinate (Q) at (4,4);
  \draw (P) -- (F2);
  \draw (F2) -- (Q);
  \draw[dashed] (Q) -- (4,0);
  \draw[dashed] (P) -- (-1,0);
  \fill (F1) circle (1.8pt);
  \fill (F2) circle (1.8pt);
  \fill (P) circle (1.8pt);
  \fill (Q) circle (1.8pt);
  \node[above right] at (P) {$P$};
  \node[right] at (Q) {$Q$};
  \node[below left] at (F1) {$F_1$};
  \node[below right] at (F2) {$F_2$};
  \node[below left] at (O) {$O$};
\end{tikzpicture}
\end{center}
\nopagebreak[4]
\end{question}
\begin{question}
\gkti{21}数列 $\{x_n\}$ 满足 $x_1=0$，$x_{n+1}=-x_n^2+x_n+c$（$n\in\mathbf N^*$）。

（1） 证明：$\{x_n\}$ 是递减数列的充分必要条件是 $c<0$；

（2） 求 $c$ 的取值范围，使 $\{x_n\}$ 是递增数列。

\end{question}